% Options for packages loaded elsewhere
\PassOptionsToPackage{unicode}{hyperref}
\PassOptionsToPackage{hyphens}{url}
\documentclass[
]{article}
\usepackage{xcolor}
\usepackage[margin=1in]{geometry}
\usepackage{amsmath,amssymb}
\setcounter{secnumdepth}{-\maxdimen} % remove section numbering
\usepackage{iftex}
\ifPDFTeX
  \usepackage[T1]{fontenc}
  \usepackage[utf8]{inputenc}
  \usepackage{textcomp} % provide euro and other symbols
\else % if luatex or xetex
  \usepackage{unicode-math} % this also loads fontspec
  \defaultfontfeatures{Scale=MatchLowercase}
  \defaultfontfeatures[\rmfamily]{Ligatures=TeX,Scale=1}
\fi
\usepackage{lmodern}
\ifPDFTeX\else
  % xetex/luatex font selection
\fi
% Use upquote if available, for straight quotes in verbatim environments
\IfFileExists{upquote.sty}{\usepackage{upquote}}{}
\IfFileExists{microtype.sty}{% use microtype if available
  \usepackage[]{microtype}
  \UseMicrotypeSet[protrusion]{basicmath} % disable protrusion for tt fonts
}{}
\makeatletter
\@ifundefined{KOMAClassName}{% if non-KOMA class
  \IfFileExists{parskip.sty}{%
    \usepackage{parskip}
  }{% else
    \setlength{\parindent}{0pt}
    \setlength{\parskip}{6pt plus 2pt minus 1pt}}
}{% if KOMA class
  \KOMAoptions{parskip=half}}
\makeatother
\usepackage{color}
\usepackage{fancyvrb}
\newcommand{\VerbBar}{|}
\newcommand{\VERB}{\Verb[commandchars=\\\{\}]}
\DefineVerbatimEnvironment{Highlighting}{Verbatim}{commandchars=\\\{\}}
% Add ',fontsize=\small' for more characters per line
\usepackage{framed}
\definecolor{shadecolor}{RGB}{248,248,248}
\newenvironment{Shaded}{\begin{snugshade}}{\end{snugshade}}
\newcommand{\AlertTok}[1]{\textcolor[rgb]{0.94,0.16,0.16}{#1}}
\newcommand{\AnnotationTok}[1]{\textcolor[rgb]{0.56,0.35,0.01}{\textbf{\textit{#1}}}}
\newcommand{\AttributeTok}[1]{\textcolor[rgb]{0.13,0.29,0.53}{#1}}
\newcommand{\BaseNTok}[1]{\textcolor[rgb]{0.00,0.00,0.81}{#1}}
\newcommand{\BuiltInTok}[1]{#1}
\newcommand{\CharTok}[1]{\textcolor[rgb]{0.31,0.60,0.02}{#1}}
\newcommand{\CommentTok}[1]{\textcolor[rgb]{0.56,0.35,0.01}{\textit{#1}}}
\newcommand{\CommentVarTok}[1]{\textcolor[rgb]{0.56,0.35,0.01}{\textbf{\textit{#1}}}}
\newcommand{\ConstantTok}[1]{\textcolor[rgb]{0.56,0.35,0.01}{#1}}
\newcommand{\ControlFlowTok}[1]{\textcolor[rgb]{0.13,0.29,0.53}{\textbf{#1}}}
\newcommand{\DataTypeTok}[1]{\textcolor[rgb]{0.13,0.29,0.53}{#1}}
\newcommand{\DecValTok}[1]{\textcolor[rgb]{0.00,0.00,0.81}{#1}}
\newcommand{\DocumentationTok}[1]{\textcolor[rgb]{0.56,0.35,0.01}{\textbf{\textit{#1}}}}
\newcommand{\ErrorTok}[1]{\textcolor[rgb]{0.64,0.00,0.00}{\textbf{#1}}}
\newcommand{\ExtensionTok}[1]{#1}
\newcommand{\FloatTok}[1]{\textcolor[rgb]{0.00,0.00,0.81}{#1}}
\newcommand{\FunctionTok}[1]{\textcolor[rgb]{0.13,0.29,0.53}{\textbf{#1}}}
\newcommand{\ImportTok}[1]{#1}
\newcommand{\InformationTok}[1]{\textcolor[rgb]{0.56,0.35,0.01}{\textbf{\textit{#1}}}}
\newcommand{\KeywordTok}[1]{\textcolor[rgb]{0.13,0.29,0.53}{\textbf{#1}}}
\newcommand{\NormalTok}[1]{#1}
\newcommand{\OperatorTok}[1]{\textcolor[rgb]{0.81,0.36,0.00}{\textbf{#1}}}
\newcommand{\OtherTok}[1]{\textcolor[rgb]{0.56,0.35,0.01}{#1}}
\newcommand{\PreprocessorTok}[1]{\textcolor[rgb]{0.56,0.35,0.01}{\textit{#1}}}
\newcommand{\RegionMarkerTok}[1]{#1}
\newcommand{\SpecialCharTok}[1]{\textcolor[rgb]{0.81,0.36,0.00}{\textbf{#1}}}
\newcommand{\SpecialStringTok}[1]{\textcolor[rgb]{0.31,0.60,0.02}{#1}}
\newcommand{\StringTok}[1]{\textcolor[rgb]{0.31,0.60,0.02}{#1}}
\newcommand{\VariableTok}[1]{\textcolor[rgb]{0.00,0.00,0.00}{#1}}
\newcommand{\VerbatimStringTok}[1]{\textcolor[rgb]{0.31,0.60,0.02}{#1}}
\newcommand{\WarningTok}[1]{\textcolor[rgb]{0.56,0.35,0.01}{\textbf{\textit{#1}}}}
\usepackage{graphicx}
\makeatletter
\newsavebox\pandoc@box
\newcommand*\pandocbounded[1]{% scales image to fit in text height/width
  \sbox\pandoc@box{#1}%
  \Gscale@div\@tempa{\textheight}{\dimexpr\ht\pandoc@box+\dp\pandoc@box\relax}%
  \Gscale@div\@tempb{\linewidth}{\wd\pandoc@box}%
  \ifdim\@tempb\p@<\@tempa\p@\let\@tempa\@tempb\fi% select the smaller of both
  \ifdim\@tempa\p@<\p@\scalebox{\@tempa}{\usebox\pandoc@box}%
  \else\usebox{\pandoc@box}%
  \fi%
}
% Set default figure placement to htbp
\def\fps@figure{htbp}
\makeatother
\setlength{\emergencystretch}{3em} % prevent overfull lines
\providecommand{\tightlist}{%
  \setlength{\itemsep}{0pt}\setlength{\parskip}{0pt}}
\usepackage{bookmark}
\IfFileExists{xurl.sty}{\usepackage{xurl}}{} % add URL line breaks if available
\urlstyle{same}
\hypersetup{
  pdftitle={Pearson and Spearman correlation},
  hidelinks,
  pdfcreator={LaTeX via pandoc}}

\title{Pearson and Spearman correlation}
\author{}
\date{\vspace{-2.5em}}

\begin{document}
\maketitle

\textbf{Inference labs} use the five-step template: Hypothesis →
Visualise → Assumptions → Conduct → Conclude.

\subsection{Learning objectives}\label{learning-objectives}

\begin{itemize}
\tightlist
\item
  Compute Pearson and Spearman correlations with \texttt{cor.test()} and
  interpret the coefficient and its CI.
\item
  Distinguish linear association from monotonic association.
\item
  Use visualisation to diagnose situations where a correlation
  coefficient misleads (Anscombe, DatasauRus).
\end{itemize}

\subsection{Prerequisites}\label{prerequisites}

Basic inferential statistics.

\subsection{Background}\label{background}

Pearson's \emph{r} measures linear association. It is the standardised
covariance of two variables; it is sensitive to outliers and assumes a
roughly linear, bivariate-normal relationship. Spearman's \emph{ρ}
measures monotonic association. It is Pearson's \emph{r} applied to the
ranks of the variables; it is robust to outliers and to non-linear
monotonic transformations.

Neither coefficient captures non-linear, non-monotonic relationships
well. Anscombe's quartet and the DatasauRus dozen are the canonical
demonstrations that a coefficient without a plot can lie --- datasets
engineered to have identical Pearson \emph{r} but radically different
shapes. The conclusion is not to abandon correlations but never to
report one without first plotting the data.

A correlation coefficient is not a causal quantity. Two variables can be
perfectly correlated without one causing the other --- a lurking
variable, a common cause, a coincidence of sample selection. The word
``correlation'' is often used loosely in the press to mean ``weak
causation''; in this course it means exactly what \texttt{cor.test()}
returns.

\subsection{Setup}\label{setup}

\begin{Shaded}
\begin{Highlighting}[]
\FunctionTok{library}\NormalTok{(tidyverse)}
\FunctionTok{library}\NormalTok{(datasauRus)}
\FunctionTok{set.seed}\NormalTok{(}\DecValTok{42}\NormalTok{)}
\FunctionTok{theme\_set}\NormalTok{(}\FunctionTok{theme\_minimal}\NormalTok{(}\AttributeTok{base\_size =} \DecValTok{12}\NormalTok{))}
\end{Highlighting}
\end{Shaded}

\subsection{1. Hypothesis}\label{hypothesis}

Two scenarios:

\begin{enumerate}
\def\labelenumi{\arabic{enumi}.}
\tightlist
\item
  A linear relationship between two simulated variables.
\item
  A monotonic but non-linear relationship where Pearson misleads.
\end{enumerate}

H0 in each case: ρ = 0. H1: ρ ≠ 0.

\subsection{2. Visualise}\label{visualise}

\begin{Shaded}
\begin{Highlighting}[]
\NormalTok{x }\OtherTok{\textless{}{-}} \FunctionTok{rnorm}\NormalTok{(n)}
\NormalTok{y\_linear   }\OtherTok{\textless{}{-}} \FloatTok{0.6} \SpecialCharTok{*}\NormalTok{ x }\SpecialCharTok{+} \FunctionTok{rnorm}\NormalTok{(n, }\DecValTok{0}\NormalTok{, }\FloatTok{0.8}\NormalTok{)}
\NormalTok{y\_monotone }\OtherTok{\textless{}{-}} \FunctionTok{exp}\NormalTok{(x) }\SpecialCharTok{+} \FunctionTok{rnorm}\NormalTok{(n, }\DecValTok{0}\NormalTok{, }\FloatTok{0.5}\NormalTok{)}
\NormalTok{df }\OtherTok{\textless{}{-}} \FunctionTok{tibble}\NormalTok{(x, y\_linear, y\_monotone) }\SpecialCharTok{|\textgreater{}}
  \FunctionTok{pivot\_longer}\NormalTok{(}\FunctionTok{starts\_with}\NormalTok{(}\StringTok{"y\_"}\NormalTok{), }\AttributeTok{names\_to =} \StringTok{"type"}\NormalTok{, }\AttributeTok{values\_to =} \StringTok{"y"}\NormalTok{)}
\end{Highlighting}
\end{Shaded}

\begin{Shaded}
\begin{Highlighting}[]
\NormalTok{df }\SpecialCharTok{|\textgreater{}}
  \FunctionTok{ggplot}\NormalTok{(}\FunctionTok{aes}\NormalTok{(x, y)) }\SpecialCharTok{+}
  \FunctionTok{geom\_point}\NormalTok{(}\AttributeTok{alpha =} \FloatTok{0.5}\NormalTok{) }\SpecialCharTok{+}
  \FunctionTok{facet\_wrap}\NormalTok{(}\SpecialCharTok{\textasciitilde{}}\NormalTok{ type, }\AttributeTok{scales =} \StringTok{"free"}\NormalTok{) }\SpecialCharTok{+}
  \FunctionTok{geom\_smooth}\NormalTok{(}\AttributeTok{method =} \StringTok{"lm"}\NormalTok{, }\AttributeTok{se =} \ConstantTok{FALSE}\NormalTok{, }\AttributeTok{colour =} \StringTok{"firebrick"}\NormalTok{, }\AttributeTok{linewidth =} \FloatTok{0.5}\NormalTok{)}
\end{Highlighting}
\end{Shaded}

\subsection{3. Assumptions}\label{assumptions}

Pearson assumes linearity and roughly bivariate-normal data. Spearman
assumes monotonicity. Neither handles non-monotonic relationships. For
both, the observations must be independent; both are sensitive to
extreme sample-size-induced significance when \emph{n} is huge.

\subsection{4. Conduct}\label{conduct}

\begin{Shaded}
\begin{Highlighting}[]
\NormalTok{p1\_spearman }\OtherTok{\textless{}{-}} \FunctionTok{cor.test}\NormalTok{(x, y\_linear,   }\AttributeTok{method =} \StringTok{"spearman"}\NormalTok{, }\AttributeTok{exact =} \ConstantTok{FALSE}\NormalTok{)}
\NormalTok{p2\_pearson  }\OtherTok{\textless{}{-}} \FunctionTok{cor.test}\NormalTok{(x, y\_monotone, }\AttributeTok{method =} \StringTok{"pearson"}\NormalTok{)}
\NormalTok{p2\_spearman }\OtherTok{\textless{}{-}} \FunctionTok{cor.test}\NormalTok{(x, y\_monotone, }\AttributeTok{method =} \StringTok{"spearman"}\NormalTok{, }\AttributeTok{exact =} \ConstantTok{FALSE}\NormalTok{)}

\FunctionTok{tibble}\NormalTok{(}
  \AttributeTok{relationship =} \FunctionTok{c}\NormalTok{(}\StringTok{"linear"}\NormalTok{, }\StringTok{"linear"}\NormalTok{, }\StringTok{"monotone"}\NormalTok{, }\StringTok{"monotone"}\NormalTok{),}
  \AttributeTok{method =} \FunctionTok{c}\NormalTok{(}\StringTok{"Pearson"}\NormalTok{, }\StringTok{"Spearman"}\NormalTok{, }\StringTok{"Pearson"}\NormalTok{, }\StringTok{"Spearman"}\NormalTok{),}
  \AttributeTok{r\_or\_rho =} \FunctionTok{c}\NormalTok{(p1\_pearson}\SpecialCharTok{$}\NormalTok{estimate, p1\_spearman}\SpecialCharTok{$}\NormalTok{estimate,}
\NormalTok{               p2\_pearson}\SpecialCharTok{$}\NormalTok{estimate, p2\_spearman}\SpecialCharTok{$}\NormalTok{estimate),}
  \AttributeTok{p =} \FunctionTok{c}\NormalTok{(p1\_pearson}\SpecialCharTok{$}\NormalTok{p.value, p1\_spearman}\SpecialCharTok{$}\NormalTok{p.value,}
\NormalTok{        p2\_pearson}\SpecialCharTok{$}\NormalTok{p.value, p2\_spearman}\SpecialCharTok{$}\NormalTok{p.value)}
\NormalTok{)}
\end{Highlighting}
\end{Shaded}

The monotonic non-linear relationship gives a much higher Spearman than
Pearson --- the rank correlation sees the perfect monotonicity.

\subsubsection{The DatasauRus
demonstration}\label{the-datasaurus-demonstration}

\begin{Shaded}
\begin{Highlighting}[]
  \FunctionTok{group\_by}\NormalTok{(dataset) }\SpecialCharTok{|\textgreater{}}
  \FunctionTok{summarise}\NormalTok{(}\AttributeTok{mean\_x =} \FunctionTok{mean}\NormalTok{(x),}
            \AttributeTok{mean\_y =} \FunctionTok{mean}\NormalTok{(y),}
            \AttributeTok{sd\_x =} \FunctionTok{sd}\NormalTok{(x),}
            \AttributeTok{sd\_y =} \FunctionTok{sd}\NormalTok{(y),}
            \AttributeTok{r =} \FunctionTok{cor}\NormalTok{(x, y), }\AttributeTok{.groups =} \StringTok{"drop"}\NormalTok{)}
\FunctionTok{head}\NormalTok{(ds\_summary)}
\end{Highlighting}
\end{Shaded}

\begin{Shaded}
\begin{Highlighting}[]
\NormalTok{datasaurus\_dozen }\SpecialCharTok{|\textgreater{}}
  \FunctionTok{ggplot}\NormalTok{(}\FunctionTok{aes}\NormalTok{(x, y)) }\SpecialCharTok{+}
  \FunctionTok{geom\_point}\NormalTok{(}\AttributeTok{alpha =} \FloatTok{0.4}\NormalTok{, }\AttributeTok{size =} \FloatTok{0.8}\NormalTok{) }\SpecialCharTok{+}
  \FunctionTok{facet\_wrap}\NormalTok{(}\SpecialCharTok{\textasciitilde{}}\NormalTok{ dataset) }\SpecialCharTok{+}
  \FunctionTok{labs}\NormalTok{(}\AttributeTok{x =} \ConstantTok{NULL}\NormalTok{, }\AttributeTok{y =} \ConstantTok{NULL}\NormalTok{)}
\end{Highlighting}
\end{Shaded}

Thirteen datasets; all have mean \emph{x} ≈ 54, mean \emph{y} ≈ 47, SDs
≈ 17 and 26, and Pearson \emph{r} ≈ −0.06. The shapes are wildly
different.

\subsection{5. Concluding statement}\label{concluding-statement}

\begin{quote}
In the linear case, Pearson \emph{r} =
\texttt{round(p1\_pearson\$estimate,\ 2)} (\emph{p} =
\texttt{signif(p1\_pearson\$p.value,\ 2)}) and Spearman ρ =
\texttt{round(p1\_spearman\$estimate,\ 2)}; the two agreed. In the
monotonic non-linear case, Pearson \emph{r} =
\texttt{round(p2\_pearson\$estimate,\ 2)} was substantially smaller than
Spearman ρ = \texttt{round(p2\_spearman\$estimate,\ 2)}, revealing that
the association was strong but not linear. The DatasauRus dozen
reproduced identical summary statistics across 13 radically different
shapes, reinforcing the rule: never report a correlation without the
scatter plot.
\end{quote}

Correlation is a single number. Its honest use requires a picture
alongside.

\subsection{Common pitfalls}\label{common-pitfalls}

\begin{itemize}
\tightlist
\item
  Reporting a Pearson \emph{r} on data with an obvious curvilinear
  shape.
\item
  Using \emph{r} to compare groups with different sample sizes; small
  \emph{r} can be ``significant'' at huge \emph{n}.
\item
  Inferring causation from a correlation.
\item
  Computing correlations on data that include outliers without a
  robustness check.
\end{itemize}

\subsection{Further reading}\label{further-reading}

\begin{itemize}
\tightlist
\item
  Anscombe FJ (1973). \emph{Graphs in Statistical Analysis}.
\item
  Matejka J, Fitzmaurice G (2017). \emph{Same stats, different graphs:
  generating datasets with varied appearance and identical statistics}.
\end{itemize}

\subsection{Session info}\label{session-info}

\begin{Shaded}
\begin{Highlighting}[]

\end{Highlighting}
\end{Shaded}

\subsection{See also --- chapter index}\label{see-also-chapter-index}

\begin{itemize}
\tightlist
\item
  \href{../index.Rmd}{Methods in this chapter}
\item
  \href{../../../ch00-find-your-method.Rmd}{Find your method (Chapter
  0)}
\end{itemize}

\end{document}
